%%%%%%%%%%%%%%%%%%%%%%%%%
%% SECTION 9: Proof of \thmref{B}
%%%%%%%%%%%%%%%%%%%%%%%%%

%%%%%%%%%%%%%%%%%%%%%%%%%%%
\section{Proof of \thmref{B}, formulas (1) and (2)}\label{sec:B}
%%%%%%%%%%%%%%%%%%%%%%%%%%%

\thmref{B} provides three formulas for Milnor's $\mu$-invariant of a three-component link $L$ in $S^3$ whose pairwise linking numbers vanish:

$(1)$\ (differential forms) \  
$\mu(L) \ = \ \displaystyle{\frac12\, \int_{T^3} \delta(\varphi*\omega_L) \wedge \omega_L}$

$(2)$\ (vector fields) \ 
$\mu(L) \ = \ \displaystyle{\frac12\, \int_{\tor\times\tor} \vec_L(\x) \times \vec_L(\y) \dtw 
\nabla_{\!\y}\,\varphi\left(\x - \y\right)\ d\x\, d\y}$

$(3)$\ (Fourier series) \ 
$\mu(L) \ = \  \displaystyle{8\pi^3 \sum_{\n\ne\zero} \a_\n\times \b_\n\dtw \n/\|\n\|^2}$

\begin{comment}
In the first integral, $\delta$  is the exterior co-derivative and  $\varphi*\omega_L$ is the convolution of  $\varphi$  and  $\omega_L$ $($taken with respect to the abelian group structure on the torus$)$.  In the second, $\nabla_{\!\y}$ indicates the gradient with respect to $\x$, the difference $\x-\x'$ is taken in the abelian group $\tor$,  and $d\x$ and $d\x'$ are volume elements.  
\end{comment}

\noindent
The notation is explained in the introduction.

In the subsections that follow, we obtain explicit expressions for the characteristic $2\textup{-form}$ $\omega_L$ and vector field $\vec_L$, and for the fundamental solution $\ph$ to the scalar Laplacian on $\tor$, and then establish formula (1) by making use of J.\,H.\,C.~Whitehead's integral formula for the Hopf invariant. Then, we show how to obtain formula (2) from formula (1). Finally, in 
\secref{B'} we review the calculus of differential forms and Fourier analysis on $T^3$, in order to make explicit the role and form of $\ph$ and obtain formula (3). 

%%
\part{Explicit formula for the 2-form $\omega_L$ and the vector field $\vec_L$ on $T^3$}
%%

Recall the formula for the characteristic map $g_L\colon T^3\to S^2$ of the link $L$ given in \propref{char},
$$
g_L(s,t,u) \ = \ \frac{F(x,y,z)}{\| F(x,y,z)\|},
$$
where $x=x(s)$, $y=y(t)$ and $z=z(u)$ parametrize the components of $L$,
and where $F$ was given in \secref{characteristic} as
$$
F(x,y,z) \ = \  (ix\dt y + iy\dt z + iz\dt x \ , \  jx\dt y + jy\dt z + jz\dt x \ , \  kx\dt y + ky\dt z + kz\dt x).
$$
Here we view $x$, $y$ and $z$ as quaternions for the multiplication and as vectors in $\reals^4$ when performing the dot product.   For simplicity, we write $g_L=F/\| F\|$, suppressing the appearance of the evaluation map $e_L(s,t,u)=(x(s),y(t),z(u))$.

Again let $\omega$ be the Euclidean area 2-form on the unit 2-sphere 
$S^2\subset\reals^3$, normalized so that the total area is 1 instead of $4\pi$.
If $\p$ is a point of $S^2$, and $\a$ and $\b$ are tangent vectors to $S^2$
at $\p$, then 
$$
\omega_{\p}(\a,\b) \ = \ \frac1{4\pi}(\a\times\b)\dtw\p.
$$
This 2-form $\omega$ on $S^2$ extends to a closed 2-form $\bah{\omega}$ in 
$\reals^3-\zero$ given by
$$
\bah{\omega}_{\p}(\a,\b) \ = \ {(\a\times\b)\dt\p\over 4\pi\|\p\|^3},
$$
which is the pullback of $\omega$ from $S^2$ to $\reals^3-\zero$ via the 
map $\p\mapsto \p/\|\p\|$.

Hence the pullback $g_L^\dast\omega$ of $\omega$ from $S^2$ to $T^3$ via
$g_L=F/\| F\|$ is the same as the pullback $F^\dast\bah{\omega}$
of $\bah{\omega}$ from $\reals^3-\zero$ to $T^3$ via $F$.

Write
$$g_L^\dast\omega \ = \ F^\dast\bah{\omega} 
\ = \ a(s,t,u)dt\wedge du+b(s,t,u)du\wedge ds+c(s,t,u)ds\wedge dt.$$
Then we have
$$
\aligned
a(s,t,u)& \ = \  F^\dast\bah\omega({\partial_t},{\partial_u}) \ = \  \bah\omega(F_\dast{\partial_t}, F_\dast{\partial_u}) \\
& \ = \ \bah{\omega}(F_t,F_u) \ = \  \frac{(F_t\times F_u)\dt F}{4\pi\| F\|^3},
\endaligned
$$
and likewise for $b(s,t,u)$ and $c(s,t,u)$, where the subscripts on $F$ denote
partial derivatives.

Therefore, the characteristic 2-form of the link $L$ is
$$\omega_L \ = \  g_L^\dast\omega \ = \ 
\frac{F_t\times F_u\dt F}{4\pi\| F\|^3}  \, dt\wedge du
+\frac{F_u\times F_s\dt F}{4\pi\| F\|^3} \, du\wedge ds
+\frac{F_s\times F_t\dt F}{4\pi\| F\|^3} \, ds\wedge dt,$$
and its corresponding characteristic vector field is 
$$\vec_L \ = \ 
\frac{F_t\times F_u\dt F}{4\pi\| F\|^3} \, \partial_s
+\frac{F_u\times F_s\dt F}{4\pi\| F\|^3} \, \partial_t
+\frac{F_s\times F_t\dt F}{4\pi\| F\|^3} \, \partial_u.$$

%%
\part{Proof of \thmref{B}, formula (1)}
%%

Let $L$ be a three-component link in the 3-sphere $S^3$ with pairwise
linking numbers $p$, $q$ and $r$ all zero.  By the first part of \thmref{A} these numbers are the degrees of the characteristic map $g_L\colon T^3\to S^2$ on the $2$-dimensional coordinate subtori. Since these degrees are all zero, $g_L$ is homotopic to a map $g\colon T^3\to S^2$ which collapses the 2-skeleton of $T^3$ to a point:
$$
g_L\ \simeq \ g\colon T^3 \stackrel{\sigma}{\longrightarrow}S^3 \stackrel{f}{\longrightarrow}S^2,
$$
where $\sigma$ is the collapsing map.  By the second part of \thmref{A}, Milnor's $\mu$-invariant of $L$ is equal to half of Pontryagin's $\nu$-invariant $\nu(g_L)$ (comparing $g_L$ to the constant map $f_{000}$ as explained in \secref{pontryagin}), which in turn is just the Hopf invariant of $f\colon S^3\to S^2$,
$$
\mu(L)\ = \ {1\over 2}\nu(g_L)\ = \ {1\over 2}\Hopf(f).
$$

We can thus use J.\,H.\,C.~Whitehead's integral formula for the Hopf invariant as the first ingredient in our formula for the $\mu$ invariant.  Starting from 
Hopf's definition of his invariant of a map $f\colon S^3\to S^2$ as the linking number between the inverse images of two regular values, Whitehead \cite{Whitehead} found an integral formula for $\Hopf(f)$ as follows.

Let $\omega$ be the area 2-form on $S^2$, normalized so that $\int_{S^2}
\omega=1$. Its pullback $f^\dast\omega$ is a closed 2-form on $S^3$,
which is exact because $H^2(S^3;\reals)=0$. Hence $f^\dast\omega=d\alpha$ for some 1-form $\alpha$ on $S^3$, and Whitehead showed that the Hopf invariant of $f$ is given by the formula
$$
\Hopf(f) \ = \ \int_{S^3}\alpha\wedge f^\dast\omega \ 
= \ \int_{S^3}\alpha\wedge d\alpha,
$$
the value of the integral being independent of the choice of $\alpha$. 

To make Whitehead's formula explicit requires a way to produce a 1-form
$\alpha$ whose differential $d\alpha$ is a given 2-form $f^\dast\omega$ on $S^3$ which is known to be exact. We give an explicit formula for this in the 
Appendix, since it is not needed in the body of the paper. Instead, we are
going to pull the whole situation back to $T^3$ and perform our calculations
there.

In particular, the formula
$$
\Hopf(f) \ = \  \int_{S^3}\alpha_0\wedge f^\dast\omega,
$$
where $\alpha_0$ is any 1-form on $S^3$ such that $d\alpha_0=f^\dast\omega$, pulls back to the formula
$$
\nu(g_L)\ = \ \int_{T^3}\bah{\alpha_0}\wedge g^\dast\omega,
$$
where $\bah{\alpha_0}$ is any 1-form on $T^3$ such that $d\bah{\alpha_0} = g^\dast\omega$.  One such choice of $\bah{\alpha_0}$ would be $\bah{\alpha_0} = \sigma^\dast\alpha_0$, with $\alpha_0$
obtained as in the Appendix.   However, since we don't have an explicit analytic formula
for $\sigma$, and since the formula for $\alpha_0$ is complicated, we will pursue 
a slightly different approach.

Taking advantage of the fact that $g_L$ is homotopic to $g$, we also have
$$
\nu(g_L) \ = \  \int_{T^3}\alpha\wedge g_L^\dast\omega \ = \ \int_{T^3}\alpha\wedge \omega_L
$$
for any $1$-form $\alpha$ on $T^3$ such that $d\alpha=\omega_L$.  Therefore,
what we really need is a canonical way to produce a $1$-form $\alpha$ on $T^3$ whose exterior derivative is the exact 2-form $\omega_L$.   We will prove in \propref{alp} that the $1$-form 
$$
\alpha_L \ = \ \delta(\ph * \omega_L)
$$
has the desired property that $d\alpha_L=\omega_L$.  Moreover, if $\widetilde\alpha$ is any other $1$-form such that $d\widetilde\alpha=\omega_L$, then $\|\alpha_L\| \le \|\widetilde\alpha\|$ in $L^2(\tor)$, with equality if and only if $\widetilde\alpha=\alpha_L$. 

We thus obtain the explicit integral formula for Milnor's $\mu$-invariant of the three component link $L$:
$$
\mu(L) \ = \ {1\over 2}\nu(g_L) \ = \ {1\over 2}\int_{T^3}\delta(\ph * \omega_L)\wedge\omega_L,
$$
which is formula (1) of \thmref{B}.

%\break

%%
\part{Proof of  \thmref{B}, formula (2)}
%%

The 1-form $\delta(\ph * \omega_L)$ on $T^3$ which appears in formula (1) of \thmref{B}
converts to the vector field $\nabla\times(\ph * \vec_L)$, which can be regarded as the magnetic field 
on $T^3$ due to the current flow $\vec_L$. The customary minus sign here is now hidden in the 
definitions of the Laplacian and its inverse (Green's operator).

The integral formula for Milnor's $\mu$-invariant given in formula (1),
$$
\mu(L) \ = \ {1\over 2}\int_{T^3}\delta(\ph * \omega_L)\wedge\omega_L\,,
$$
then converts to the formula
$$
\mu(L) \ = \ {1\over 2}\int_{T^3}\left(\vphantom{x^2}\nabla\times(\ph * \vec_L)\right)\dt\vec_L\,d\vol
$$
in the language of vector fields. To obtain formula (2) of \thmref{B}, we first expand out the convolution integral
$$
(\ph * \vec_L)(\x) \ = \ \int_{T^3}\vec_L(\y)\,\ph(\x-\y)\,d\y.
$$
Then we compute its curl:
$$
\begin{aligned}
\left(\vphantom{x^2}\nabla_{\!\x}\times(\ph * \vec_L)\right)(\x)
\ &= \  \int_{T^3}\nabla_{\!\x}\times\left(\vphantom{x^2}\vec_L(\y)\,\ph(\x-\y)\right)\,d\y  \\
\ &= \  -\int_{T^3}\vec_L(\y)\times\nabla_{\!\x}\,\ph(\x-\y)\,d\y\,,
\end{aligned}
$$
using the product formula $\nabla\times(f\veca)=f(\nabla\times \veca)-\veca\times\nabla f$.
Note that $\nabla_{\!\x}\times\vec_L(\y)=\zero$ since, from the point of view of the variable
$\x$, the vector field $\vec_L(\y)$ is constant.

Inserting the expression for the curl of $\ph*\vec_L$ into the formula for $\mu(L)$, we get
$$
\aligned
\mu(L)
&\ = \ -\frac12\int_{T^3\times T^3}\left(\vphantom{x^2}\vec_L(\y)\times\nabla_{\!\x}\,\ph(\x-\y)\right)\dtw
\vec_L(\x)\,d\x\,d\y \\
&\ = \ -\frac12\int_{T^3\times T^3}\vec_L(\x)\times\vec_L(\y)\dtw\nabla_{\!\x}\,\ph(\x-\y)\,d\x\,d\y\\
&\ = \ \frac12\int_{T^3\times T^3}\vec_L(\x)\times\vec_L(\y)\dtw\nabla_{\!\y}\,\ph(\x-\y)\,d\x\,d\y\,,
\endaligned
$$
where at the last step we hid the minus sign by taking the gradient of $\ph$ with respect to
$\y$ instead of $\x$, and obtained
formula (2) of \thmref{B}.


%%%%%%%%%%%%%%%%%%%%%%%%%
%% SECTION 9: Proof of \thmref{B}, formula (3)
%%%%%%%%%%%%%%%%%%%%%%%%%

%%%%%%%%%%%%%%%%%%%%%%%%%%%
\section{Fourier series and the Proof of \thmref{B}, formula (3)}\label{sec:B'}
%%%%%%%%%%%%%%%%%%%%%%%%%%%


\vskip -.2in
\vskip -.2in
%%
\part{Fourier series and the fundamental solution of the Laplacian}
%%

In the proofs of formulas (1) and (2) of \thmref{B}, we needed to find a 1-form $\alpha$ on $T^3$ whose exterior derivative is the exact 2-form $\omega_L$ associated with the three-component link $L$ in $S^3$. We asserted that we can choose  $\alpha=\delta(\ph * \omega_L)$, where $\ph$ is the fundamental solution of the scalar Laplacian on $T^3$. Furthermore, we asserted that this choice of $\alpha$ is canonical in the sense that it has the smallest $L^2$ norm among all possible choices. We justify these assertions in this section and lay the groundwork for the proof of formula (3) of \thmref{B} by studying the calculus of differential forms on $T^3=(\br/2\pi\bz)^3  $ in terms of their Fourier series. 

We will prove the following two results:

%%%
\begin{comment}
\noindent\bf Theorem C\rm.\ \ \it
\rm (a) \it The fundamental solution 
of the scalar Laplacian on the $3$-torus \\
$T^3=S^1 \times S^1 \times S^1$ is given by the formula
$$
\varphi(x,y,z) \ = \  
{1\over 8\pi^3}\sum_{{ m,n,p = -\infty }\atop
{m^2+n^2+p^2\neq 0}}^\infty 
{e^{i(mx+ny+pz)}\over m^2+n^2+p^2}.
$$
The function $\varphi$ is $C^\infty$ at all points of $T^3$ except $(0,0,0)$, where it becomes infinite. 
\rm (b) \it If $\omega$ is any exact differential form on $T^3$ with $C^\infty$ coefficients, then 
$$
\alpha \ = \ \delta(\ph * \omega)
$$
satisfies $d\alpha=\omega$. Furthermore, if $d\widetilde\alpha=\omega$ as well, then $\Vert\alpha\Vert_{L^2}
\le\Vert\widetilde\alpha\Vert_{L^2}$, with equality if and only if $\widetilde\alpha=\alpha$.
\rm 
\end{comment}
%%%

\begin{proposition}\label{prop:fs}
\textbf{\boldmath  \ \ The fundamental solution of the scalar Laplacian on the\\
 $3$-torus $T^3 = (\br/2\pi\bz)^3$ is given by the formula
$$
\varphi(\x) \ = \
{1\over 8\pi^3}\sum_{\n\ne\zero}
e^{i\n\dte\x} / \|\n\|^2.
$$
The function $\varphi$ is $C^\infty$ at all points $
\x\in T^3$ except $\zero$, where it becomes infinite.} 
\end{proposition}

\bigskip

\begin{proposition}\label{prop:alp}
\textbf{\boldmath  If $\omega$ is any exact differential form on $T^3$ with $C^\infty$ coefficients, then 
$$
\alpha \ = \ \delta(\ph * \omega)
$$
is a $C^\infty$ differential form satisfying $d\alpha=\omega$.  Furthermore, if $d\widetilde\alpha=\omega$ as well, then $\Vert\alpha\Vert_{L^2}\le\Vert\widetilde\alpha\Vert_{L^2}$, with equality if and only if $\widetilde\alpha=\alpha$.} 
\end{proposition} 


Before diving into calculations, we pause for some words of explanation: We write $\Omega^k(T^3)$ for the space of $C^\infty$ $k$-forms on $T^3$. With $d$ as the exterior differentiation operator taking $\Omega^k(T^3)$ to $\Omega^{k+1}(T^3)$, and $\delta$ the co-differentiation map adjoint to  $d$ in the $L^2$ sense, the Laplacian of a $k$-form $\alpha$ is 
$$
\Delta\alpha \ = \ (d\delta+\delta d)\alpha.
$$
This definition gives us the ``geometer's sign convention'' for the Laplacian on functions (0-forms):
$$
\Delta f \ = \  -\left({\partial^2 f\over\partial x_1^2}+{\partial^2 f\over\partial x_2^2} +
{\partial^2 f\over\partial x_3^2}\right).
$$


%%
\part{Proof of \propref{fs}}
%%

The \it fundamental solution \rm of the scalar Laplacian is a function $\varphi$, 
convolution with which ``inverts'' the Laplacian to the extent that this is possible. On $T^3$, only 
functions that integrate to zero are in the range of the Laplacian, and so ``the'' fundamental 
solution of $\Delta$ on $T^3$ is the function $\varphi$ which satisfies
$$
\int_{T^3}\varphi\,d\!\vol \ = \  0
\qquad\textup{and}\qquad 
\Delta(\varphi * f) \ = \ f
$$
for all $f\in C^\infty(T^3)$ such that $\displaystyle{\int_{T^3}f\,d\!\vol=0}$.

Even though we have expressed $\varphi$ in terms of complex exponentials, 
the value of $\varphi$ is real for real values of $\x$ because of the symmetry of the coefficients.

\figref{fundsol} shows the graph of the corresponding fundamental solution
$$\varphi(\x) \ = \
{1\over 4\pi^2}\sum_{\n\ne\zero} e^{i\n\dte \x} /  \|\n\|^2 
$$
of the scalar Laplacian on the 2-torus $S^1 \times S^1$, summed for $|\n| \leq 10$, and displayed over the range $[-3\pi,3\pi]\times[-3\pi,3\pi]$.

%%% FIGURE 31: Fund Sol %%%
\begin{figure}[h!]
\includegraphics[height=260pt]{FigFundSolDennis.pdf}
\caption{Fundamental solution of the scalar Laplacian on $S^1 \times S^1$}
\label{fig:fundsol}
\end{figure}
%%%%%%%%%%%%%%%%%%  


With respect to the $L^2$ inner product
$$
\langle f,g\rangle \ = \  \int_{T^3}f(\x)\,\overline{g(\x)}\,d\x\,,
$$
the set of complex exponentials $\{e^{i\n\dte\x} \st \n\in\bz^3\}$ is an orthogonal set, and for all $\n$, 
we have $\langle e^{i\n\dte\x},e^{i\n\dte\x}\rangle = 8\pi^3$. 

Any function $f$ in $L^2(T^3)$ can be expanded into a Fourier series
$$
f\ \sim \ \sum_{\n\in\ints^3} c_{\n}e^{i\n\dte \x}$$
where 
$$
c_\n\ = \ {1\over 8\pi^3}\langle f,e^{i\n\dte\x}\rangle \ = \  {1\over 8\pi^3}\int_{T^3} f(\x)\,e^{-i\n\dte\x}\,d\x.
$$
Moreover, by Plancharel's theorem, the $L^2(T^3)$-norm of $f$ is equal to the $\ell^2$-norm of the 
sequence of coefficients:
$$
\Vert f\Vert_{L^2} \ = \  \int_{T^3}|f(x)|^2\,d\x \ = \  \sum_{\n\in\ints^3} 8\pi^3|c_\n|^2.
$$
This shows that the function $\varphi$ of \propref{fs} is in $L^2(T^3)$, since $\|c_\n\|^2$ is of order $1/\|\n\|^4$, while the number of lattice points at distance $\|\n\|$
from $\zero$ is of order $\|\n\|^2$. 

Because the range of the scalar Laplacian is the (closed) subspace of functions $f$ for which 
$$
\int_{T^3} f(\x)\,d\x\ = \ 0\,,
$$
we will denote this subspace of $L^2(T^3)$ by $L^2_0(T^3)$, and similarly for other function spaces, where a zero subscript will indicate that all functions in the space have average value zero.  In terms of Fourier coefficients, $f\in L^2_0(T^3)$ if and only if the Fourier coefficient $c_\zero=c_{(0,0,0)}=0$.  

For functions $f$ and $g$ in $C^\infty(T^3)$, their {\itb convolution}
$f*g$ is defined by
$$
\left(\vphantom{x^2}f*g\right)(\x)\ = \ \int_{T^3}f(\y)\,g(\x-\y)\,d\y\,,
$$
and also lies in $C^\infty(T^3)$.  One checks easily that $
e^{i\m\dte\x}*e^{i\n\dte\x} = 0$ if $\m\neq \n$, and 
$$
e^{i\n\dte\x}*e^{i\n\dte\x}\ = \ 8\pi^3 e^{i\n\dte\x}.
$$
Therefore, if $f$ and $g$ have Fourier series $\sum c_{\n}e^{i\n\dte\x}$ and
$\sum d_{\n}e^{i\n\dte\x}$ respectively, then the Fourier series of $f*g$ is
$$
8\pi^3\sum_{\n\in\ints^3}c_{\n}d_{\n}e^{i\n\dte\x}.
$$

From this
it follows that the operation of convolution satisfies $f*g=g*f$
and $(f*g)*h=f*(g*h)$ whenever all the convolutions are defined. Furthermore,
$$
{\partial\over \partial x_k}(f*g)\ = \ {\partial f\over \partial x_k}*g \ = \  f*{\partial g\over \partial x_k}.
$$

For $u\in C^\infty(T^3)$, with 
$$
u(\x)\ = \ \sum_{\n\in\ints^3} b_{\n}e^{i\n\dte\x}\,,
$$
we have
$$
\Delta u(\x)\ = \ \sum_{\n\neq\zero}|\n|^2b_{\n}e^{i\n\dte\x}.
$$
Because the $(0,0,0)$-coefficient of $\Delta u$ is zero, this shows why in order to solve $\Delta u=f$ we must restrict $f$ to be in $C^\infty_0(T^3)$.   And for $f\in C^\infty_0(T^3)$, with
$$
f(\x)\ = \ \sum_{\n\neq\zero}c_{\n}e^{i\n\dte\x},
$$
the unique solution $u\in C^\infty_0(T^3)$ of $\Delta u=f$ has $b_{\n}= c_\n/|\n|^2$,
and so, since 
$$
\ph(\x)\ = \ \frac{1}{8\pi^3}\sum_{\n\ne\zero}{e^{i\n\dte \x}}/{\|\n\|^2}\,,
$$
we conclude that $u=\varphi * f$.  This completes the proof of \propref{fs} except for the assertion about the smoothness of $\ph$ away from its singularity, which follows from well-known local regularity theorems for elliptic partial differential equations (see, for example, Folland~\cite{Folland},  Theorem 6.33).

%%
\part{Fourier series and the calculus of differential forms on the 3-torus}
%%

To prove \propref{alp}, we express the calculus of differential forms on $T^3$ in terms of Fourier series.  Except for the fundamental solution $\ph$ of the (scalar) Laplacian on $T^3$, we will assume all functions and forms are $C^\infty$.

We continue to express functions (0-forms) as Fourier series: for 
$\x=(s,t,u)$ and $\n=(n_1,n_2,n_3)$, we write
$$
f(\x) \ = \ \sum_{\n\in\ints^3}c_{\n}e^{i\n\dte\x}\ \in\ \Omega^0(T^3).
$$
Likewise, we employ the notation described in the introduction and express 1-forms using 
$\c_{\n}=(c_\n^{s},c_\n^{t},c_\n^{u})$
and $d\x=(ds,dt,du)$. We write
$$
\alpha(\x) \ = \ \sum_{\n\in\ints^3}\c_{\n}e^{i\n\dte\x}\dt d\x\ \in\ \Omega^1(T^3).
$$
With $\star d\x=(dt\wedge du,du\wedge ds,ds\wedge dt)$ as before, we can express a 2-form as
$$
\beta(\x) \ = \ \sum_{\n\in\ints^3}\c_{\n}e^{i\n\dte\x}\dt \star d\x\ \in\ \Omega^2(T^3).
$$
Finally, with $dV=ds\wedge dt\wedge du$, a 3-form can be written as
$$
\gamma(\x) \ = \ \sum_{n\in\ints^3}c_{\n}e^{i\n\dte\x}dV\ \in\ \Omega^3(T^3).
$$

It is straightforward to express both the exterior derivative $d$ and the
codifferential $\delta$ in terms of Fourier coefficients. With $f$, $\alpha$,
$\beta$ and $\gamma$ as above, we have
$$
\begin{array}{lll}
\displaystyle{df \ = \ \sum ic_{\n}\n e^{i\n\dte\x}\dt d\x} &
\mbox{\rm\ and\ } & \delta f \ = \ 0,\\[0.3cm]
\displaystyle{d\alpha \ = \ \sum i\n\times \c_\n e^{i\n\dte\x}\dt \star d\x} &
\mbox{\rm\ and\ } & \displaystyle{\delta\alpha \ = \ -\sum i\n\dt \c_\n e^{i\n\dte\x}},\\[0.3cm]
\displaystyle{d\beta \ = \ \sum i\n\dt\c_\n e^{i\n\dte\x}dV} &
\mbox{\rm\ and\ } & \displaystyle{\delta\beta \ = \ \sum i\n\times\c_\n e^{i\n\dte\x}
\dt d\x,}\\[0.3cm]
d\gamma \ = \ 0 & \mbox{\rm\ and\ } & 
\displaystyle{\delta\gamma \ = \ -\sum ic_\n \n e^{i\n\dte\x}\dt\star d\x}.
\end{array}
$$
Thus, exterior differentiation and co-differentiation are expressed in terms of vector algebraic operations on the Fourier coefficients.
From these expressions, we conclude the following about the kernel and image
of $d$ and $\delta$:

\begin{itemize}
\item For 0-forms, 
$$
\begin{aligned}
\ker d &\ = \ \{f \st c_\n=0\ \mbox{\rm for}\ \n\ne\zero,
\mbox{\rm\ while\ }
\ c_\zero
\ \mbox{\rm is arbitrary}\} \ \ \ \\[0.2cm]
\im d\, &\ = \  \{0\} \\[0.2cm]
\ker \delta &\ = \  \{\mbox{\rm all}\ f\} \\[0.2cm]
\im \delta\, &\ = \  \{f \st c_\zero=0\}
\end{aligned}
$$
\item For 1-forms,
$$
\begin{aligned}
\ker d &\ = \  \{\alpha \st \c_\n=\lambda_\n\n\ \mbox{\rm for}\ \n\ne\zero,
\mbox{\rm\ while\ }\c_\zero
\ \mbox{\rm is arbitrary}\} \\[0.2cm] 
\im d\, &\ = \  \{\alpha \st \c_\n=\lambda_\n\n\ \mbox{\rm for}\ \n\ne\zero,\ \c_\zero=\zero\} \\[0.2cm]
\ker \delta &\ = \  \{\alpha \st \c_\n\dt\n =0\ \mbox{\rm for}\ \n\ne\zero,
\mbox{\rm\ while\ }\c_\zero \ \mbox{\rm is arbitrary}\} \\[0cm]
%%%%%%% attempt to even out the spacing
\im \delta\, &\ = \  \{\alpha \st \c_\n\dt\n=0\ \mbox{\rm for}
\ \n\ne\zero,\ \c_\zero=\zero\}
\end{aligned}
$$
\item For 2-forms,
$$
\begin{aligned}
\ker d &\ = \  \{\beta \st \c_\n\dt\n=0\ \mbox{\rm for}\ \n\ne\zero,
\mbox{\rm\ while\ }\c_\zero
\ \mbox{\rm is arbitrary}\} \\[0.01cm]
\im d\, &\ = \  \{\beta \st \c_\n\dt\n=0\ \mbox{\rm for}\ \n\ne\zero,\ \c_\zero=\zero\} \\[0cm]
\ker \delta &\ = \  \{\beta \st \c_\n=\lambda_\n\n\ \mbox{\rm for}\ \n\ne\zero,
\mbox{\rm\ while\ }\c_\zero \ \mbox{\rm is arbitrary}\} \\[0cm]
\im \delta\, &\ = \  \{\beta \st \c_\n=\lambda_\n\n\ \mbox{\rm for}
\ \n\ne\zero,\ \c_\zero=\zero\}
\end{aligned}
$$
\item For 3-forms,
$$
\begin{aligned}
\ker d&\ = \ \{\mbox{\rm all}\ \gamma\}\\[0.15cm]
\im d\, &\ = \  \{\gamma \st c_\zero=0\} \\[0.2cm]
\ker \delta &\ = \  \{\gamma \st c_\n=0\ \mbox{\rm for}\ \n\ne\zero,
\mbox{\rm\ while\ }\ c_\zero
\ \mbox{\rm is arbitrary}\} \\[0.1cm]
\im \delta\, &\ = \ \{0\}
\end{aligned}
$$
\end{itemize}

Therefore $\im d\subset\ker d$ and $\im \delta\subset\ker\delta$, and
the following orthogonal (with respect to the $L^2$ inner product) decompositions hold:
$$
\Omega^k(T^3)\ = \ \ker\delta\oplus\im d \ = \ \im\delta\oplus\ker d
$$
for $k=0,\ldots,3$. The $k$-forms in $\ker d\cap\ker\delta$ are the forms
whose Fourier series contain only constant terms. These are called \it harmonic \rm
$k$-forms because they are in the kernel of
the Laplacian $\Delta=d\delta+\delta d$, so we write $\Har^k(T^3)=\ker d\cap\ker\delta$.
We then have the Hodge decomposition 
$$
\Omega^k(T^3)\ = \ \im\delta\oplus\Har^k(T^3)\oplus \im d.
$$
The Laplacian $\Delta\colon\Omega^k(T^3)\to\Omega^k(T^3)$ preserves this Hodge decomposition, 
taking $\im \delta$ bijectively to itself, killing $\Har^k(T^3)$, and taking $\im d$ bijectively to itself.

%%%%%
\part{Proof of \propref{alp}}
%%%%%

We must show that if $\omega$ is any $C^\infty$ exact differential $k$-form on $T^3$, then the $(k-1)$-form $\alpha=\delta(\ph*\omega)$ is $C^\infty$ and satisfies
$d\alpha=\omega$, and that if $d\widetilde\alpha=\omega$ as well, then $\Vert\alpha\Vert_{L^2}
\le\Vert\widetilde\alpha\Vert_{L^2}$, with equality if and only if $\widetilde\alpha=\alpha$.

Because it is the case we use in \thmref{B}, we will carry this out 
specifically for 2-forms. The proofs for 1-forms and 3-forms
are essentially the same. 

For a 2-form 
$$
\beta\ = \ \sum_{\n\in\ints^3}\c_\n e^{i\n\dte\x}\dt \star d\x\,,
$$ 
its Laplacian is given by
$$
\Delta\beta\ = \ \sum_{\n\ne\zero}|\n|^2\c_\n e^{i\n\dte\x}\dt \star d\x.
$$
We can use convolution with the fundamental solution $\ph$ of the scalar Laplacian to express the Green's operator,
$$
\Gr(\beta)\ = \ \ph * \beta \ = \  \sum_{\n\ne \zero}{\c_\n\over|\n|^2}e^{i\n\dte\x}\dt \star d\x.
$$
Clearly the Laplacian $\Delta$ and the Green's operator are inverses of one another when applied to 2-forms $\beta$ with $\c_\zero=\zero$, equivalently, to 2-forms $\beta$ orthogonal to $\Har^2(T^3)$.

To prove \propref{alp}, assume that $\beta$ is exact, that is, 
$\beta\in\im d$. Then $\Gr(\beta)$ is also exact, hence certainly closed, and therefore
$$
\beta\ = \ \Delta\Gr(\beta)\ = \ (d\delta+\delta d)\Gr(\beta)\ = \ d\delta\Gr(\beta)\ = \ d\delta(\ph*\beta).
$$
Thus $\alpha=\delta(\ph*\beta)$ satisfies $d\alpha=\beta$.

If $\widetilde\alpha$ is any other 2-form satisfying $d\widetilde\alpha=\beta$, then $\widetilde\alpha$ differs from $\alpha$ by some closed 2-form, which must be $L^2$-orthogonal to $\alpha$ since $\alpha\in\im\delta$. Thus, $\Vert\alpha\Vert_{L^2}\le\Vert\widetilde\alpha\Vert_{L^2}$
by the Pythagorean theorem.

Since the Fourier coefficients of a $C^\infty$ 2-form decrease
faster than any negative power of $|\n|$, we have that $\Gr(\beta)$ and 
$\delta\Gr(\beta)$ will be $C^\infty$ if $\beta$ is. This completes the proof
of \propref{alp}.

%%
\part{Proof of \thmref{B}, formula (3)}
%%

This formula expresses Milnor's invariant $\mu(L)$ in terms of the Fourier coefficients of $\omega_L$:
$$
\mu(L)\ = \ 8\pi^3\sum_{\n\ne\zero}\frac{\a_\n\times\b_\n\dtew\n}{\|\n\|^2}
$$
where 
$$
\omega_L\ = \ \sum_{\n\ne\zero}\c_\n e^{i\n\dte\x}\dtew \star d\x
$$
with $\c_\n=\a_\n+i\b_\n$ and with $\a_\n$ and $\b_\n$ real.

The formula for $\omega_L$ is summed over $\n\ne\zero$ because the hypothesis
of pairwise linking numbers zero is equivalent to the vanishing of the harmonic component of $\omega_L$, which in turn is equivalent to the vanishing of its Fourier coefficient $\c_\zero$.

We begin with formula (1) of \thmref{B} for $\mu(L)$:
$$
\mu(L)\ = \ {1\over 2}\int_{T^3}\delta(\ph*\omega_L)\wedge\omega_L.
$$
We have
$$
\ph*\omega_L\ = \ \sum_{\n\ne\zero}\frac{\c_\n e^{i\n\dte\x}}{\|\n\|^2}\dtew \star d\x\,,
$$
and hence, from our table of derivatives and co-derivatives in terms of Fourier
series,
$$
\delta(\ph*\omega_L)\ = \ \sum_{\n\ne\zero} \frac{(\n\times\c_\n)\,i\,e^{i\n\dte\x}}{\|\n\|^2}\dtew d\x.
$$
To compute $\delta(\ph*\omega_L)\wedge\omega_L$, we use the fact that if $\vv \dtew d\x$ is a 1-form and $\w\dtew \! \star \!d\x$ is a 2-form, then
$$
(\vv\dtew d\x)\wedge(\w\dtew \! \star\!d\x)\ = \ (\vv\dtew \w)\,dV.
$$
Thus
$$
\delta(\ph*\omega_L)\wedge\omega_L\ = \ \sum_{\m\ne\zero}\sum_{\n\ne\zero}
\frac{(\n\times\c_\n)\dtew\c_\m\,ie^{i(\m+\n)\dte\x}}{\|\n\|^2}\,dV.
$$
When we insert this double sum into the integral 
$\mu(L)={1\over 2}\int_{T^3}\delta(\ph*\omega_L)\wedge\omega_L,$
most of the terms in the summation will integrate to zero, leaving only the terms
where $\m+\n=\zero$. In those cases, $e^{i(\m+\n)\dte\x}=e^0=1$ integrates
to $8\pi^3$, the volume of $T^3$.

Thus
$$
\begin{aligned}
\mu(L)&\ = \ 4\pi^3\sum_{\n\ne\zero}\frac{i(\n\times\c_\n)\dtew\c_{-\n}}{\|\n\|^2} \\
&\ = \ 4\pi^3\sum_{\n\ne\zero}\frac{i(\c_\n\times\c_{-\n})\dtew\n}{\|\n\|^2}
\end{aligned}
$$
and it remains to simplify this last series.

Since $\c_\n=\a_\n+i\b_\n$ and since $\omega_L$ is real-valued, we have
$
\c_{-\n} = \bah{\c_\n} = \a_\n-i\b_\n.
$
Hence
$$
\c_\n\times\c_{-\n}\ = \ (\a_\n+i\b_\n)\times(\a_\n-i\b_\n)\ = \ 
-2i\,\a_\n\times\b_\n,$$
and therefore
$$
\begin{aligned}
\mu(L)&\ = \ 4\pi^3\sum_{\n\ne\zero}\frac{i(\c_\n\times\c_{-\n})\dtew\n}{\|\n\|^2}\\
&\ = \ 8\pi^3\sum_{\n\ne\zero}\frac{(\a_\n\times\b_\n)\dtew\n}{\|\n\|^2}\,,
\end{aligned}
$$
completing the proof of formula (3) and, with it, that of \thmref{B}.

%%
\part{Numerical computation}
%%

We used Maple and Matlab 
to calculate an approximation to Milnor's $\mu$-invariant for the three-component link $L$
in $S^3$ given by
$$
\aligned
x(s)&\ = \ [\textstyle{4\over 5}\sin s,\cos s,\textstyle{3\over 5}\sin s,0], \\
y(t)&\ = \ [\textstyle{4\over 5}\sin t,0,\cos t,\textstyle{3\over 5}\sin t], \\
z(u)&\ = \ [\textstyle{4\over 5}\sin u,\textstyle{3\over 5}\sin u,0,\cos u],
\endaligned
$$
for $s\in[0,2\pi]$,
$t\in[0,2\pi]$ and $u\in[0,2\pi]$,
which is a concrete realization of the Borromean rings with $\mu=-1$. 

%%%%%%%%%%%%%
%%% FIGURE 32: Borromean Rings %%%
\begin{figure}[h!]
\includegraphics[height=130pt]{FigBorRafal.pdf}
\caption{Borromean rings}
\label{fig:bor}
\end{figure}
%%%%%%%%%%%%%%%%%%  
%%%%%%%%%%%%%

Maple gave the (long, ugly) formula for $\omega_L$ obtained from 
\propref{char}, and we used Matlab to calculate approximations to its Fourier coefficients $\c_\n$
for divisions of the $s$, $t$ and $u$ intervals into 60 subintervals and for $\n\in[-15,15]^3$. The approximation of $\mu$ we obtained in this way was $-0.97$. 

%%%%
\appendix
%%%%

\section*{Appendix.  Whitehead's integral formula for the Hopf invariant}

We take this opportunity to make J.\,H.\,C.~Whitehead's integral formula for the Hopf invariant of a map $f\colon S^3\to S^2$ explicit, and have relegated it to an appendix because we do not use the result in this paper.

Recall from \secref{B} Whitehead's formula, 
$$
\Hopf(f) \ = \ \int_{S^3}\alpha\wedge\omega_f,
$$
where $\omega_f=f^\dast\omega$ is the pullback to $S^3$ of the normalized area form $\omega$ on $S^2$, and $\alpha$ is any 1-form on $S^3$ such that $d\alpha=\omega_f$. To make this formula explicit requires a way to produce such an $\alpha$. 

To do this, we first write Whitehead's formula in the language of vector fields,
$$
\Hopf(f) \ = \ \int_{S^3}\veca\dtw\vec_f\ d\vol,
$$
where the 2-form $\omega_f$ on $S^3$ has been converted to the vector field $\vec_f$ in the usual way, and where $\veca$ is any vector field on $S^3$ such that $\nabla\times\veca=\vec_f$.

An explicit recipe for $\veca$ was given by DeTurck and Gluck~\cite{DeTurckGluck}:
$$
\veca(\y) \ = \ \BS(\vec_f)(\y) \ = \ \int_{S^3}P_{\y\x}\vec_f(\x)\times\nabla_{\!\y}\,\ph(\x,\y)\,d\x.
$$
Here $\BS$ is the Biot--Savart operator for vector fields on the $3$-sphere.  In the last integral, $P_{\y\x}$ indicates parallel transport in $S^3$ along the geodesic segment from $\x$ to $\y$, and the function $\ph$ is given by
$$
\ph(\alpha) \ = \ -{1\over 4\pi^2}(\pi-\alpha)\csc\alpha,
$$
and $\ph(\x,\y)$ is an abbreviation for $\ph(\alpha(\x,\y))$, with
$\alpha(\x,\y)$ the geodesic distance on $S^3$ between $\x$ and $\y$. 

The significance of the above function $\ph$ is that it is the fundamental solution of a shifted 
Laplacian on $S^3$: 
$$-\Delta\ph-\ph=\delta,$$
where $\delta$ is the Dirac delta function. 

The above formula for $\BS(\vec_f)$ is the analogue on $S^3$ of the classical Biot--Savart formula from 
electrodynamics in $\reals^3$, expressing the magnetic field $\BS(\vec_f)$ in terms of the current 
flow $\vec_f$. The equation $\nabla\times\BS(\vec_f)=\vec_f$ is just one of Maxwell's 
equations, transplanted to $S^3$.

Inserting this formula for $\BS(\vec)$ into the previous formula for the 
Hopf invariant and performing simple manipulations, we get
$$
\Hopf(f) \ = \ -\int_{S^3\times S^3} P_{\y\x}\vec_f(\x)\times\vec_f(\y)\dtw
\nabla_{\!\y}\,\ph(\x,\y)\,d\x\,d\y,
$$
the explicit version of Whitehead's integral formula on the 3-sphere $S^3$.



